% Options for packages loaded elsewhere
\PassOptionsToPackage{unicode}{hyperref}
\PassOptionsToPackage{hyphens}{url}
%
\documentclass[
  10pt,
  letterpaper,
]{article}
\usepackage{amsmath,amssymb}
\usepackage{iftex}
\ifPDFTeX
  \usepackage[T1]{fontenc}
  \usepackage[utf8]{inputenc}
  \usepackage{textcomp} % provide euro and other symbols
\else % if luatex or xetex
  \usepackage{unicode-math} % this also loads fontspec
  \defaultfontfeatures{Scale=MatchLowercase}
  \defaultfontfeatures[\rmfamily]{Ligatures=TeX,Scale=1}
\fi
\usepackage{lmodern}
\ifPDFTeX\else
  % xetex/luatex font selection
\fi
% Use upquote if available, for straight quotes in verbatim environments
\IfFileExists{upquote.sty}{\usepackage{upquote}}{}
\IfFileExists{microtype.sty}{% use microtype if available
  \usepackage[]{microtype}
  \UseMicrotypeSet[protrusion]{basicmath} % disable protrusion for tt fonts
}{}
\makeatletter
\@ifundefined{KOMAClassName}{% if non-KOMA class
  \IfFileExists{parskip.sty}{%
    \usepackage{parskip}
  }{% else
    \setlength{\parindent}{0pt}
    \setlength{\parskip}{6pt plus 2pt minus 1pt}}
}{% if KOMA class
  \KOMAoptions{parskip=half}}
\makeatother
\usepackage{xcolor}
\usepackage[margin=1in]{geometry}
\usepackage{color}
\usepackage{fancyvrb}
\newcommand{\VerbBar}{|}
\newcommand{\VERB}{\Verb[commandchars=\\\{\}]}
\DefineVerbatimEnvironment{Highlighting}{Verbatim}{commandchars=\\\{\}}
% Add ',fontsize=\small' for more characters per line
\newenvironment{Shaded}{}{}
\newcommand{\AlertTok}[1]{\textcolor[rgb]{1.00,0.00,0.00}{\textbf{#1}}}
\newcommand{\AnnotationTok}[1]{\textcolor[rgb]{0.38,0.63,0.69}{\textbf{\textit{#1}}}}
\newcommand{\AttributeTok}[1]{\textcolor[rgb]{0.49,0.56,0.16}{#1}}
\newcommand{\BaseNTok}[1]{\textcolor[rgb]{0.25,0.63,0.44}{#1}}
\newcommand{\BuiltInTok}[1]{\textcolor[rgb]{0.00,0.50,0.00}{#1}}
\newcommand{\CharTok}[1]{\textcolor[rgb]{0.25,0.44,0.63}{#1}}
\newcommand{\CommentTok}[1]{\textcolor[rgb]{0.38,0.63,0.69}{\textit{#1}}}
\newcommand{\CommentVarTok}[1]{\textcolor[rgb]{0.38,0.63,0.69}{\textbf{\textit{#1}}}}
\newcommand{\ConstantTok}[1]{\textcolor[rgb]{0.53,0.00,0.00}{#1}}
\newcommand{\ControlFlowTok}[1]{\textcolor[rgb]{0.00,0.44,0.13}{\textbf{#1}}}
\newcommand{\DataTypeTok}[1]{\textcolor[rgb]{0.56,0.13,0.00}{#1}}
\newcommand{\DecValTok}[1]{\textcolor[rgb]{0.25,0.63,0.44}{#1}}
\newcommand{\DocumentationTok}[1]{\textcolor[rgb]{0.73,0.13,0.13}{\textit{#1}}}
\newcommand{\ErrorTok}[1]{\textcolor[rgb]{1.00,0.00,0.00}{\textbf{#1}}}
\newcommand{\ExtensionTok}[1]{#1}
\newcommand{\FloatTok}[1]{\textcolor[rgb]{0.25,0.63,0.44}{#1}}
\newcommand{\FunctionTok}[1]{\textcolor[rgb]{0.02,0.16,0.49}{#1}}
\newcommand{\ImportTok}[1]{\textcolor[rgb]{0.00,0.50,0.00}{\textbf{#1}}}
\newcommand{\InformationTok}[1]{\textcolor[rgb]{0.38,0.63,0.69}{\textbf{\textit{#1}}}}
\newcommand{\KeywordTok}[1]{\textcolor[rgb]{0.00,0.44,0.13}{\textbf{#1}}}
\newcommand{\NormalTok}[1]{#1}
\newcommand{\OperatorTok}[1]{\textcolor[rgb]{0.40,0.40,0.40}{#1}}
\newcommand{\OtherTok}[1]{\textcolor[rgb]{0.00,0.44,0.13}{#1}}
\newcommand{\PreprocessorTok}[1]{\textcolor[rgb]{0.74,0.48,0.00}{#1}}
\newcommand{\RegionMarkerTok}[1]{#1}
\newcommand{\SpecialCharTok}[1]{\textcolor[rgb]{0.25,0.44,0.63}{#1}}
\newcommand{\SpecialStringTok}[1]{\textcolor[rgb]{0.73,0.40,0.53}{#1}}
\newcommand{\StringTok}[1]{\textcolor[rgb]{0.25,0.44,0.63}{#1}}
\newcommand{\VariableTok}[1]{\textcolor[rgb]{0.10,0.09,0.49}{#1}}
\newcommand{\VerbatimStringTok}[1]{\textcolor[rgb]{0.25,0.44,0.63}{#1}}
\newcommand{\WarningTok}[1]{\textcolor[rgb]{0.38,0.63,0.69}{\textbf{\textit{#1}}}}
\usepackage{longtable,booktabs,array}
\usepackage{calc} % for calculating minipage widths
% Correct order of tables after \paragraph or \subparagraph
\usepackage{etoolbox}
\makeatletter
\patchcmd\longtable{\par}{\if@noskipsec\mbox{}\fi\par}{}{}
\makeatother
% Allow footnotes in longtable head/foot
\IfFileExists{footnotehyper.sty}{\usepackage{footnotehyper}}{\usepackage{footnote}}
\makesavenoteenv{longtable}
\setlength{\emergencystretch}{3em} % prevent overfull lines
\providecommand{\tightlist}{%
  \setlength{\itemsep}{0pt}\setlength{\parskip}{0pt}}
\setcounter{secnumdepth}{-\maxdimen} % remove section numbering
\usepackage{xurl}
\setlength{\emergencystretch}{3em}
\usepackage{microtype}
\ifLuaTeX
  \usepackage{selnolig}  % disable illegal ligatures
\fi
\usepackage{bookmark}
\IfFileExists{xurl.sty}{\usepackage{xurl}}{} % add URL line breaks if available
\urlstyle{same}
\hypersetup{
  hidelinks,
  pdfcreator={LaTeX via pandoc}}

\author{}
\date{}

\begin{document}

\section{DDTA BA2 semantic vocabulary candidate -
R1}\label{ddta-ba2-semantic-vocabulary-candidate---r1}

\textbf{Status:} PROVISIONAL CANDIDATE / NOT CLOSED / BA2 OPEN

\textbf{Derived by:} BA2-T2

\textbf{Repository baseline reviewed:}
\texttt{f87d05e5ea1cee969246e5eae1dd73b8b6c3a5a1}

\textbf{Structural dependency:}
\texttt{BA2\_PROPOSITION\_STRUCTURE\_CANDIDATE\_R1.md}

\textbf{Closed identity dependency:}
\texttt{BAReferent\ +\ BAProposition}

\subsection{1. Purpose and boundary}\label{purpose-and-boundary}

This artifact records the smallest vocabulary architecture that survived
BA2-T2 pressure testing. It is not an exhaustive verb dictionary and
does not close BA2.

The semantic contract remains:

\begin{Shaded}
\begin{Highlighting}[]
\NormalTok{BAProposition}
\NormalTok{|{-} semanticOperatorKey              1}
\NormalTok{|{-} participation                    1..*}
\NormalTok{|    |{-} roleKey                     1}
\NormalTok{|    \textasciigrave{}{-} term                        1}
\NormalTok{|{-} polarity                         1  [provisional normalized modifier; default affirmative]}
\NormalTok{\textasciigrave{}{-} scopedModifier                   0..*}
\end{Highlighting}
\end{Shaded}

Every accepted semantic operator key must be defined by a registry-level
semantic contract rather than by its display word alone.

At minimum, an operator contract must make mechanically available:

\begin{Shaded}
\begin{Highlighting}[]
\NormalTok{operator key}
\NormalTok{normative method{-}neutral meaning}
\NormalTok{compatible participation roles}
\NormalTok{required/optional role cardinalities}
\NormalTok{modifier compatibility or exclusions}
\end{Highlighting}
\end{Shaded}

Human labels and authoring synonyms are downstream lexical concerns.
They must not become the semantic identity of an operator merely because
a document uses a particular verb.

\subsection{2. Operator key versus authoring
predicate}\label{operator-key-versus-authoring-predicate}

The two reviewed corpora use many surface predicates such as
\texttt{creates}, \texttt{records}, \texttt{produces}, \texttt{reads},
\texttt{sends}, \texttt{requestsCapture}, \texttt{references},
\texttt{confirmsHandoff}, \texttt{updates}, \texttt{increases} and
\texttt{decreasesReserved}.

BA2-T2 rejects mechanical copying of that list into the Base Analysis
vocabulary.

Different surface verbs can share one project-semantic relation, while
similar words can carry different semantics in context. The canonical
unit is therefore a \textbf{stable semantic operator key with a
normative definition and role contract}, not an authoring token.

The separation is:

\begin{Shaded}
\begin{Highlighting}[]
\NormalTok{governed/source lexical predicate}
\NormalTok{        {-}\textgreater{} derivation mapping in BA3}
\NormalTok{stable BA2 semantic operator key}
\NormalTok{        {-}\textgreater{} canonical assertion semantics}
\NormalTok{readable canonical label / synonym}
\NormalTok{        {-}\textgreater{} BA5 lexical assistance}
\end{Highlighting}
\end{Shaded}

BA2-T2 defines only the middle semantic contract.

\subsection{3. Provisional operator-family
facet}\label{provisional-operator-family-facet}

A small family facet is useful for registry organization and
compatibility checks, but BA2-T2 does not treat it as a new identity
family or as sufficient assertion semantics.

The current candidate families are:

\begin{itemize}
\tightlist
\item
  \texttt{ACTION} - an occurrence, operation, observation, creation,
  transfer or state-affecting behavior is asserted;
\item
  \texttt{RELATION} - a reusable association, dependency, correlation,
  realization, responsibility or consumption relationship is asserted;
\item
  \texttt{CONSTRAINT} - admissible behavior/state/acceptance or
  applicability is restricted;
\item
  \texttt{CLASSIFICATION} - a reusable method-neutral semantic kind is
  assigned to a referent.
\end{itemize}

\textbf{Status:}
\texttt{PROVISIONAL\ ORGANIZING\ FACET\ /\ NOT\ A\ CLOSED\ ENUMERATION}.

A consumer must select by the exact operator key when exact semantics
matter. \texttt{ACTION} or \texttt{RELATION} alone is too coarse.

\subsection{4. Seed semantic operator
keys}\label{seed-semantic-operator-keys}

The table below is a falsifiable \textbf{seed set}, not a complete
vocabulary. Keys are written in readable form for the candidate; later
BA2 work may separate machine-stable identifiers from canonical labels.

\begin{longtable}[]{@{}
  >{\raggedright\arraybackslash}p{(\columnwidth - 6\tabcolsep) * \real{0.2500}}
  >{\raggedright\arraybackslash}p{(\columnwidth - 6\tabcolsep) * \real{0.2500}}
  >{\raggedright\arraybackslash}p{(\columnwidth - 6\tabcolsep) * \real{0.2500}}
  >{\raggedright\arraybackslash}p{(\columnwidth - 6\tabcolsep) * \real{0.2500}}@{}}
\toprule\noalign{}
\begin{minipage}[b]{\linewidth}\raggedright
Candidate key
\end{minipage} & \begin{minipage}[b]{\linewidth}\raggedright
Family
\end{minipage} & \begin{minipage}[b]{\linewidth}\raggedright
Minimum semantic distinction preserved
\end{minipage} & \begin{minipage}[b]{\linewidth}\raggedright
Representative pressure
\end{minipage} \\
\midrule\noalign{}
\endhead
\bottomrule\noalign{}
\endlastfoot
\texttt{transfer} & ACTION & content is conveyed from a source to a
destination & facial-access FR-3.4 \\
\texttt{produce} & ACTION & an actor/capability makes an output/result
available & availability, reservation, payment, carrier results \\
\texttt{create} & ACTION & a new project-semantic item/event occurrence
is established & OrderEvaluation, Reservation, FulfillmentTask,
StockMovement \\
\texttt{observe} & ACTION & an actor/capability reads/queries/observes
existing project state without asserting state creation & availability
read, provider status resolution \\
\texttt{transition} & ACTION & a subject changes from/to governed state
or lifecycle disposition & Order, Reservation, FulfillmentTask/result
lifecycle \\
\texttt{correlate} & RELATION & one item/result/operation is bound to
the same evaluation/request/context identity & RecognitionRequest,
OrderEvaluation, provider operations \\
\texttt{reference} & RELATION & one referent explicitly points to
another without implying correlation, dependency or ownership &
FulfillmentTask -\textgreater{} Order/Reservation \\
\texttt{dependOn} & RELATION & one project meaning requires another
capability/result/milestone as prerequisite & MR/service and handoff
dependencies \\
\texttt{consumeService} & RELATION & a consumer uses a
capability/service without transferring ownership & cross-MR service
consumption, external transport service \\
\texttt{realize} & RELATION & a more concrete realization
implements/materializes an abstract project meaning & connectivity
-\textgreater{} Ethernet/Wi-Fi \\
\texttt{assignResponsibility} & RELATION & responsibility/authority is
placed on a party/capability for a scope & camera vs processor, internal
inventory authority \\
\texttt{ownOrManage} & RELATION & ownership/management authority over a
scoped resource/service is asserted or denied & external transport
ownership boundary \\
\texttt{constrain} & CONSTRAINT & admissible behavior, state, result,
acceptance or applicability is restricted & completion, confidentiality,
integrity, provenance, all-or-nothing policy \\
\texttt{classify} & CLASSIFICATION & a referent is assigned a reusable
method-neutral semantic kind & information/resource, behavior/event,
capability, contract, state/context \\
\end{longtable}

\subsubsection{Seed-set discipline}\label{seed-set-discipline}

The presence of a document verb does not automatically add an operator
key. A new operator key is justified only if collapsing it into an
existing key would erase a semantic distinction required by at least one
governed corpus or method-neutral consumer.

Conversely, two keys should merge if regression shows their distinction
adds no mechanically relevant project semantics.

\subsection{5. Operator distinctions that BA2-T2 refuses to
collapse}\label{operator-distinctions-that-ba2-t2-refuses-to-collapse}

Current evidence requires at least these distinctions to remain
mechanically visible:

\begin{Shaded}
\begin{Highlighting}[]
\NormalTok{correlation           != generic reference}
\NormalTok{service consumption   != ownership/management transfer}
\NormalTok{provider raw state    != governed project state}
\NormalTok{request acceptance    != physical handoff}
\NormalTok{realization           != generic descriptive qualifier}
\NormalTok{referent kind         != participation role}
\NormalTok{project constraint    != method{-}owned threat interpretation}
\end{Highlighting}
\end{Shaded}

These are semantic boundaries, not claims about how many classes or
database tables an implementation uses.

\subsection{6. Participation-role
contract}\label{participation-role-contract}

BA2-T2 rejects one unconstrained global role list whose labels are
assumed to mean the same thing under every operator.

Roles have reusable semantic meanings, but \textbf{validity and
cardinality are operator-scoped}.

A role registry entry may therefore be reusable, while an operator
contract decides whether that role is required, optional, repeatable or
forbidden.

Provisional reusable role concepts include:

\begin{longtable}[]{@{}
  >{\raggedright\arraybackslash}p{(\columnwidth - 2\tabcolsep) * \real{0.5000}}
  >{\raggedright\arraybackslash}p{(\columnwidth - 2\tabcolsep) * \real{0.5000}}@{}}
\toprule\noalign{}
\begin{minipage}[b]{\linewidth}\raggedright
Role concept
\end{minipage} & \begin{minipage}[b]{\linewidth}\raggedright
Method-neutral meaning
\end{minipage} \\
\midrule\noalign{}
\endhead
\bottomrule\noalign{}
\endlastfoot
\texttt{actor} & party/capability that performs or causes an action \\
\texttt{source} & origin endpoint/location of a transfer \\
\texttt{destination} & receiving endpoint/location of a transfer \\
\texttt{content} & information/resource conveyed or acted upon \\
\texttt{request} & request/command/input that initiates or scopes an
interaction \\
\texttt{result} & governed output/result produced by an action \\
\texttt{context} & referent providing
correlation/applicability/observation context \\
\texttt{consumer} & party/capability consuming a service/capability \\
\texttt{provider} & party/capability providing the consumed service \\
\texttt{service} & consumed capability/service meaning \\
\texttt{dependent} & meaning whose validity/operation depends on
another \\
\texttt{prerequisite} & required dependency target \\
\texttt{subject} & referent whose state/classification/constraint is
asserted \\
\texttt{fromState} & previous state when materially governed and
known \\
\texttt{toState} & resulting state when materially governed and known \\
\texttt{abstract} & abstract project meaning in a realization
relation \\
\texttt{realization} & concrete realization of the abstract meaning \\
\texttt{responsibleParty} & party/capability assigned
responsibility/authority \\
\texttt{responsibilityScope} & project meaning over which responsibility
is assigned \\
\texttt{classifiedReferent} & referent being method-neutrally
classified \\
\texttt{semanticKind} & controlled semantic-kind value used by
classification \\
\texttt{constraintTarget} & referent whose admissible semantics are
constrained \\
\texttt{constraintValue} & controlled/local value or referent expressing
the restriction \\
\end{longtable}

\textbf{Status:}
\texttt{PROVISIONAL\ ROLE\ CONCEPTS\ /\ NOT\ CLOSED\ GLOBAL\ VOCABULARY}.

\subsubsection{Operator-scoped examples}\label{operator-scoped-examples}

A \texttt{transfer} proposition can require one \texttt{source}, one
\texttt{destination} and one \texttt{content}, while allowing zero or
more context/request terms when the governed semantics need them.

A \texttt{transition} proposition can require one \texttt{subject} and
one \texttt{toState}; \texttt{fromState} is required only when the
source semantics distinguish it.

A \texttt{consumeService} proposition can require \texttt{consumer} and
\texttt{service}, and may identify \texttt{provider} when the provider
itself is governed project meaning.

A \texttt{classify} proposition requires one \texttt{classifiedReferent}
and at least one controlled \texttt{semanticKind} value.

The exact full compatibility/cardinality matrix remains a BA2 regression
target rather than being silently treated as closed here.

\subsection{7. Why roles are
operator-scoped}\label{why-roles-are-operator-scoped}

The same word can be unsafe as a globally context-free semantic slot.

For example, \texttt{source} in a transfer means an interaction origin.
In BA3, a source locator means provenance of an analytical statement.
Those responsibilities must not collide merely because English uses the
same word.

Similarly, \texttt{target} could mean a constraint target, transition
subject, destination, or analysis-method target. Operator-scoped
contracts avoid making one ambiguous global role carry incompatible
meanings.

Thus:

\begin{Shaded}
\begin{Highlighting}[]
\NormalTok{role semantic concept may be reusable}
\NormalTok{        +}
\NormalTok{role admissibility/cardinality is operator{-}scoped}
\NormalTok{        +}
\NormalTok{lexical label is not the role\textquotesingle{}s authority}
\end{Highlighting}
\end{Shaded}

\subsection{8. Scoped-modifier contract}\label{scoped-modifier-contract}

BA2-T1 required a scoped modifier capability. BA2-T2 narrows when it may
be used.

A modifier is appropriate only when all of the following hold:

\begin{enumerate}
\def\labelenumi{\arabic{enumi}.}
\tightlist
\item
  it changes the interpretation of exactly one proposition or one
  participation binding;
\item
  it does not introduce a new independently reusable participant set;
\item
  it does not need independent assertion-level provenance/review/change
  identity;
\item
  it is not reused as project meaning by other propositions.
\end{enumerate}

If any condition fails, the semantics should be promoted to a
\texttt{BAReferent} and/or a separate \texttt{BAProposition} rather than
hidden in a modifier mini-language.

\subsubsection{Provisional modifier
kinds}\label{provisional-modifier-kinds}

The following modifier capabilities survive current pressure testing:

\begin{itemize}
\tightlist
\item
  \texttt{polarity} - affirmative/negative assertion without inventing
  separate \texttt{notX} operators;
\item
  \texttt{condition} - local precondition/guard for proposition
  applicability;
\item
  \texttt{applicability} - local applicability restriction when it has
  no independent identity;
\item
  \texttt{quantification} - counts/cardinality such as exactly-one or
  all-lines when locally scoped;
\item
  \texttt{temporalScope} - before/after/during/until style scope when
  local rather than independently reusable;
\item
  \texttt{completionOutcome} - success/failure/completeness semantics
  local to the proposition;
\item
  \texttt{atomicity} - local all-or-nothing semantic qualifier;
\item
  \texttt{concurrency} - local simultaneous/competing-operation
  qualifier;
\item
  \texttt{idempotency} - repeated handling must not repeat the semantic
  effect.
\end{itemize}

These are \textbf{candidate modifier kinds}, not a closed vocabulary.

\subsection{9. Polarity refinement}\label{polarity-refinement}

BA2-T2 strongly prefers explicit proposition polarity over proliferating
negative operator keys.

For example:

\begin{Shaded}
\begin{Highlighting}[]
\NormalTok{operator: ownOrManage}
\NormalTok{polarity: NEGATIVE}
\NormalTok{participants:}
\NormalTok{  responsibleParty {-}\textgreater{} Project}
\NormalTok{  responsibilityScope {-}\textgreater{} underlyingTransportInfrastructure}
\end{Highlighting}
\end{Shaded}

is semantically cleaner than inventing \texttt{doesNotOwnOrManage} as a
separate operator.

Likewise, \texttt{MUST\ NOT\ count\ as\ physical\ handoff} should be
represented through a constraint/polarity structure rather than a
lexical \texttt{notHandoff} operator.

\textbf{Disposition:} \texttt{EXPLICIT\ POLARITY\ -\ STRONG\ CANDIDATE}.

\subsection{10. Modifier promotion rule}\label{modifier-promotion-rule}

Several high-value semantics must \textbf{not} be collapsed into generic
modifiers by default:

\begin{itemize}
\tightlist
\item
  correlation with a reusable request/evaluation;
\item
  realization by Ethernet/Wi-Fi or another independently meaningful
  realization;
\item
  responsibility/authority placement;
\item
  service consumption/provider relationship;
\item
  reusable lifecycle state/milestone identity;
\item
  a separately governed constraint with its own provenance/change
  lifecycle.
\end{itemize}

These semantics either have reusable participants or need independent
analytical assertion identity. They therefore remain explicit
participation/relation/proposition semantics.

This promotion rule is necessary so \texttt{scopedModifier} does not
become a second untyped prose field.

\subsection{11. Classification as a
BAProposition}\label{classification-as-a-baproposition}

BA2-T2 supports classification through the same accepted assertion
mechanism rather than adding a new BA1 family or an opaque authoritative
\texttt{type} field.

Conceptually:

\begin{Shaded}
\begin{Highlighting}[]
\NormalTok{operator: classify}
\NormalTok{participants:}
\NormalTok{  classifiedReferent {-}\textgreater{} RecognitionCapture}
\NormalTok{  semanticKind       {-}\textgreater{} information{-}resource}
\end{Highlighting}
\end{Shaded}

\texttt{semanticKind} may be a controlled vocabulary value when the kind
itself does not require BAReferent identity.

This representation has three advantages:

\begin{enumerate}
\def\labelenumi{\arabic{enumi}.}
\tightlist
\item
  classification remains methodology-neutral and queryable;
\item
  BA3 can later attach source/provenance/review semantics to the
  classification assertion;
\item
  changing a classification does not mutate the identity of the
  underlying BAReferent.
\end{enumerate}

The exact semantic-kind vocabulary remains \textbf{OPEN}. Current
recurring kinds such as participant/capability, information/resource,
behavior/event, contract, store, boundary and state/context are evidence
inputs, not yet a closed taxonomy.

\subsection{12. Facial-access instantiation
pressure}\label{facial-access-instantiation-pressure}

The bounded FR-3.4 transfer can be represented without DFD/STRIDE
categories:

\begin{Shaded}
\begin{Highlighting}[]
\NormalTok{operator: transfer}
\NormalTok{participants:}
\NormalTok{  source      {-}\textgreater{} CameraSubsystem}
\NormalTok{  destination {-}\textgreater{} RecognitionProcessor}
\NormalTok{  content     {-}\textgreater{} RecognitionCapture}
\NormalTok{  context     {-}\textgreater{} RecognitionRequest}
\NormalTok{modifiers:}
\NormalTok{  completionOutcome {-}\textgreater{} incomplete != successful}
\end{Highlighting}
\end{Shaded}

Additional propositions can preserve:

\begin{Shaded}
\begin{Highlighting}[]
\NormalTok{correlate(RecognitionCapture, RecognitionRequest)}
\NormalTok{consumeService(projectDelivery, localConnectivity)}
\NormalTok{realize(localConnectivity, Ethernet)}
\NormalTok{NEGATIVE ownOrManage(Project, underlyingTransportInfrastructure)}
\NormalTok{constrain(deliveryBehavior, confidentiality/integrity/authorized{-}origin semantics)}
\end{Highlighting}
\end{Shaded}

A STRIDE consumer may map this project semantics to its own
transfer-oriented categories, but Base Analysis does not contain
\texttt{Process}, \texttt{DataFlow}, \texttt{TrustBoundary},
\texttt{Spoofing} or other method-owned roots.

\subsection{13. Order-fulfillment instantiation
pressure}\label{order-fulfillment-instantiation-pressure}

The order corpus exercises different operator contracts:

\begin{itemize}
\tightlist
\item
  \texttt{create} for OrderEvaluation, Reservation, StockMovement and
  FulfillmentTask identity-bearing project meanings;
\item
  \texttt{produce} for AvailabilityResult, ReservationResult,
  PaymentAuthorizationResult, PaymentSettlementResult and
  DispatchAcceptance;
\item
  \texttt{correlate} for those results/operations with OrderEvaluation,
  Order or FulfillmentTask context;
\item
  \texttt{transition} for Order/Reservation lifecycle semantics;
\item
  \texttt{consumeService} for cross-MR capability consumption without
  co-ownership;
\item
  \texttt{assignResponsibility} / \texttt{ownOrManage} for internal
  inventory authority versus external provider responsibility;
\item
  \texttt{constrain} plus local atomicity/concurrency/idempotency
  modifiers for all-or-nothing reservation and exactly-once effects;
\item
  \texttt{dependOn} for the physical handoff milestone reused by
  inventory stock issue and payment capture.
\end{itemize}

Critically, carrier request acceptance is not allowed to collapse into
the physical handoff milestone. Their distinct semantics require
distinct project referents/propositions even though both occur in one
operational chain.

\subsection{14. Rejected alternatives}\label{rejected-alternatives}

BA2-T2 rejects:

\begin{itemize}
\tightlist
\item
  copying every FunctionalPredicate/document verb into the canonical BA
  vocabulary;
\item
  using a human-readable label as the semantic identity of an operator
  or role;
\item
  one global context-free role list with no operator compatibility
  rules;
\item
  encoding negative semantics through separate \texttt{notX} operators
  by default;
\item
  putting correlation, realization, responsibility or service
  consumption into a generic modifier bag;
\item
  inferring reusable referent classification only from the roles it
  happens to occupy;
\item
  importing DFD/STRIDE element or threat-category names into the shared
  vocabulary;
\item
  promoting operator keys, role keys or modifier kinds into new BA1 BAE
  identity families.
\end{itemize}

\subsection{15. What remains open}\label{what-remains-open}

BA2 is not closed because important questions still need regression:

\begin{itemize}
\tightlist
\item
  whether all seed operator keys remain necessary across both corpora or
  some should merge/split;
\item
  the complete operator-role compatibility/cardinality matrix;
\item
  machine-stable key format versus canonical display labels and aliases;
\item
  whether the provisional family facet adds value or can collapse
  without loss;
\item
  edge cases for modifier promotion, especially nested conditions and
  concurrency semantics;
\item
  the minimal seed set of referent semantic kinds needed for
  projections;
\item
  whether a later holdout forces a new method-neutral semantic
  distinction.
\end{itemize}

\subsection{16. Falsification rule}\label{falsification-rule}

Revise this candidate if regression shows any of the following:

\begin{enumerate}
\def\labelenumi{\arabic{enumi}.}
\tightlist
\item
  a seed operator distinction can collapse without loss of mechanically
  required project semantics;
\item
  one semantic distinction cannot be expressed without raw-prose
  reconstruction under the proposed operator/role/modifier contract;
\item
  operator-scoped role contracts prevent a legitimate cross-operator
  reusable role meaning that requires a different mechanism;
\item
  a modifier must routinely carry its own participants/provenance/reuse,
  showing that it should instead be a proposition/referent;
\item
  classification-as-proposition cannot support method-neutral projection
  without creating cyclic or ambiguous semantics;
\item
  the bounded method consumer requires a method-owned category in the
  shared core;
\item
  a genuine third identity family is required, triggering the BA1 reopen
  criteria rather than being hidden inside BA2.
\end{enumerate}

\subsection{17. Next microstep}\label{next-microstep}

Execute only:

\begin{quote}
\textbf{BA2-T3 - cross-corpus regression of the operator/role/modifier
candidate and semantic-key/lexical separation.}
\end{quote}

BA2 remains open. BA3 remains not started.

\end{document}
